\honest{A Sense That More Is Possible}{Eliezer Yudkowsky}{2009}

People tell me they are not interested in rationality because they have ``known a couple of rational people and they didn't seem any happier.'' Who were these people? ``Probably an Objectivist or some such,'' I guess, or an ordinary scientist, or an ordinary atheist. My answer: ``That's really not a whole lot of rationality.'' Nor is being able to derive Bayes's Theorem, which would eliminate perhaps 98\% of them; it is a basic theorem. I have always had a sense that there ought to be an art of thinking whose students become visibly more competent and formidable. I do not see it. I see a hint of it in a few very senior psychologists, such as Robyn Dawes, Daniel Gilbert, and Tooby and Cosmides, who care so much about rationality that it probably makes their colleagues uncomfortable. Even that is not a lot. I doubt that the mastery of those who impress me compares with ``John Conway's mastery of math.'' The knowledge we actually use is probably less than a nuclear engineer's, perhaps less than a construction engineer's knowledge of bridges, and our self-taught practice is less than an Olympic runner's, or maybe a professional tennis player's. The root, I suspect, is that we have never got together and systematized our skills; each of us built them alone. \nb{So no rational person the objector could name would count. Later in the post the question changes from whether rational people are happier to why rationalists are not ``surrounded by a visible aura of formidability.''}

Why are rationalists not surrounded by a visible aura of formidability, or found at the top of every elite selected for anything to do with thought? Why do most seem like ordinary people of somewhat above-average intelligence with one more hobbyhorse? One answer is that they get less systematic training ``than a first-dan black belt gets in hitting people.'' I include myself; I am no beisutsukai, because one person can create only so much of an art without statistics on results. I know a single use of rationality, ``reduction of confusing cognitions.'' A mature training program would teach other arts that would make me stronger and happier, but my life has room for only one sub-art built from scratch. In that one skill I can ``punch through brick'' and am ``working on steel along my way to adamantine,'' but I have only ``a mere casual street-fighter's grasp of how to kick or throw or block.''

The chief obstacle is that it is hard to test whether a training program works. There are experiments on debiasing now and again, ``but it tends to be something like, `Make the students practice this for an hour, then test them two weeks later.''' I want something more like running half the applicants to a three-month summer program through version A and half through version B, and surveying them five years later. I cite none of the existing experiments. \nb{Darrin Lehman and Richard Nisbett had followed students through four years of university, and reported in 1990 that ``reasoning is taught and generalizable.'' It was not a randomized comparison of programs, which is what the post asks for, but it measured years of training, not an hour.}

Why are there schools of martial arts and no rationality dojos? Not because hitting matters more than thinking. It is easier to verify that you have hit someone. And people want to hit, and believe that a systematic art can make them visibly more formidable. Long ago some people had the sense that more was possible; they shared, practiced and formalized their techniques because they thought they should be awesome and were willing to work for it. Martial arts schools got somewhere ``because they could tell when they had hit someone; and the schools competed against each other regularly in realistic contests with clearly-defined winners.''

Why is there no art of rationality? I give three reasons and number them backwards. ``Third,'' rationalists have trouble working in groups. ``Second,'' success in training is hard to verify. ``But first,'' people lack the sense that rationality should be trained like a martial art. That sense is the title of the post: that rationality should have as much knowledge behind it as nuclear engineering, superstars who practice as hard as chess grandmasters, and successful practitioners with an evident aura of awesome. Nor do people look at the lack of such formidability and say, ``We must be doing something wrong.'' Rationality seems one more hobby to talk about at parties, with few real consequences.

I quote a reader, Daniel Burfoot, who suggests why intelligence seems to matter so much for rationality: when everything is improvised with little training, intelligence is the main factor left. I blame the lack of training. \nb{The post does not consider that, apart from intelligence and knowledge and practice in particular fields, there may be little general skill of thinking to train. The lack of visibly formidable rationalists fits that explanation too.}
